\documentclass[12pt]{article}
\usepackage[utf8]{inputenc}
\usepackage[T1]{fontenc}
\usepackage{hyperref}
\title{Teribe Language Resources for K-12 Curriculum Development}
\author{Compiled by Claude Sonnet 5 (claude-sonnet-5) for Andr\'as Kornai}
\date{}
\begin{document}
\maketitle
\section{Introduction}
Teribe, also called Naso (ISO 639-3 tfr), is a Chibchan language of northwestern Panama. A Costa Rican variety is known as T\'erraba or Bror\'an. One search source puts fluent speakers at about 500 (2018), against the roughly 3,300 in the reference description. Nearly all speakers also speak Spanish.
\section{Primary Educational Resources}
Spanish has been the language of instruction in Naso schools since the 1970s. A 2009 Panamanian law created a bilingual education programme for the Naso territory. One source estimates it will take at least 10 to 15 years to train enough Naso teachers. No primers or textbooks turned up in search.
\section{Linguistic and Reference Materials}
\begin{itemize}
\item Juan Diego Quesada, \emph{A Grammar of Teribe} (Lincom, 2000), reviewed at \url{https://linguistlist.org/issues/14/1824/} (the review link was not opened, so the attribution to Quesada comes from the search summary only)
\item Overview articles: \url{https://en.wikipedia.org/wiki/Teribe_language} and \url{https://en.wikipedia.org/wiki/Naso_language}
\item Natalia Berm\'udez, Naso documentation report (2010--2011): \url{https://gbs.uni-koeln.de/sites/gbs/user_upload/Berichte_D/Bermudez_Naso_Report.pdf}
\item Costa Rican variety: \url{https://dipalicori.ucr.ac.cr/en/lengua-cultura/naso-terraba/}
\end{itemize}
\section{Religious and Community Resources}
No scripture or community organisation was identified in the results.
\section{Digital and Audio Resources}
Endangered Languages Alliance page on Naso: \url{https://elalliance.org/languages/naso/}. No audio collection or dictionary was found.
\section{Curriculum Development Context}
The legal basis for bilingual education exists, but teachers and materials are scarce. Quesada's grammar and the Berm\'udez documentation are the starting points.
\end{document}